\section{Limitations}\label{sec:limits}

Each item narrows what the results support.

\begin{enumerate}
\item \emph{Emulation, not physics.} Every upset is a bit flipped in a
  read buffer by software. Nothing here is a radiation measurement, and
  no result transfers to a beam without one.
\item \emph{The correction limit is not exercised.} At most two
  erroneous symbols reached a codeword in one decode, against eight
  correctable. Multi-bit upsets, bursts, and the regime beyond eight
  symbols, where a Reed-Solomon decoder can miscorrect silently
  (\texttt{known-limitations.md}, 3.11), are untested. The measured
  volumes did not carry \texttt{DATA\_CSUM}, which exists to catch
  that case.
\item \emph{Grouped parity.} Nine consecutive bytes in the parity of one
  codeword defeat it (\cref{sec:capsule}); burst tolerance holds for the
  3\,824 coded bytes only.
\item\label{lim:indirect} \emph{Indirect blocks.} In Reed-Solomon mode the indirect parity
  detects and journals but does not repair the buffer; pointers 478 to
  511 have no parity; a pointer flipped within the data area onto an
  allocated block is not detected without \texttt{DATA\_SELFID}
  (\cref{sec:indirect-design}). The project's own notes advise against
  Reed-Solomon as the default indirect-parity mode until a per-block
  generation covers indirect blocks (\texttt{known-limitations.md},
  3.13).
\item \emph{The instrument.} emufi 0.8.1 logs one hook a bio early on
  this kernel (\cref{sec:injector}). Below $10^{6}$\,ppm a lone flip's
  block is known only to within one block. The effective target
  included the block after each target block. emufi seeds its generator
  at load, so an attack is reproducible in its parameters, not bit for
  bit.
\item \emph{Unequal doses.} Flips follow bios, and filesystems issue
  different numbers of bios for the same read. Results are reported per
  flip received, not at equal exposure.
\item \emph{Coverage of the campaign.} Three runs; one target file per
  attack; static workload. squashfs was never exercised. In runs 2 and 3
  part of the cluster was not exercised (\cref{sec:setup}).
\item \emph{Functional coverage.} xfstests ran its \texttt{auto} group
  and declined \xfsXNotrun{} of its \xfsTotal{} tests on x86-64 and
  \xfsANotrun{} on aarch64 (\cref{tab:xfs}); what they test is untested
  here. The tree checker of the x86-64 kernel was not built into the
  aarch64 kernel (\cref{sec:setup}).
\item \emph{generic/476 on aarch64.} Within the sweep the harness
  stopped it at its time budget while it was still writing. Its reading
  of the volume images taken then, mounted and in use, found block
  pointers beyond the end of the device; an image of a volume in use
  cannot tell damage on disk from writes in flight, and the question is
  left open. Run alone on the same image the test passed, and it passed
  in the x86-64 sweep.
\item \emph{Tool reporting.} The bench and the xfstests harness
  misreported in ways listed in \cref{sec:selfaudit}; the results here
  are read from the raw records, not from their summaries.
\item \emph{Kernel log.} Correction lines are rate-limited: the
  per-decode figures rest on the first ten lines of each attack.
\item \emph{Platform.} Virtual machines on one host; virtual disks and
  USB flash media with their own controllers and error correction
  below the filesystem. Nothing was run on target hardware.
\item \emph{Write path.} emufi's write hook was armed but recorded no
  flip on \bfs; upsets on the write side are untested.
\item \emph{Size.} At \moduleLines{} lines the module is beyond the
  5\,000-line audit budget the project once set itself.
\end{enumerate}
